\section{All Tasks}
\label{app:tasks}

\begin{table}[H]
\centering
\small
\begin{tabular}{llL{10.2cm}}
\toprule
Task & Kind & What the agent must do \\
\midrule
\multicolumn{3}{l}{\emph{Train tasks}} \\
\texttt{fix\_fib}         & coding & Fix the bug in \texttt{fib.py} so the tests pass, without changing the tests. \\
\texttt{palindrome}       & coding & Write \texttt{is\_palindrome(s)}, ignoring case, spaces and punctuation, so the tests pass. \\
\texttt{fix\_import}      & coding & \texttt{python main.py} crashes; make it run and print the result, without changing \texttt{mathutils.py}. \\
\texttt{word\_count}      & coding & Write a script \texttt{count\_words.py FILE} that prints only the number of words in the file. \\
\texttt{safe\_divide}     & coding & Make \texttt{safe\_divide} return \texttt{None} when dividing by 0, so the tests pass. \\
\texttt{rename\_function} & coding & Rename the function \texttt{calc} to \texttt{calculate\_total} everywhere it is used. \\
\texttt{json\_port}       & coding & Change the port in \texttt{config.json} to the number 8080, keeping all other settings and valid JSON. \\
\texttt{organize}         & PC & Move every \texttt{.txt} file into a folder \texttt{txt} and every \texttt{.csv} file into a folder \texttt{csv}. \\
\texttt{count\_lines}     & PC & Write the number of lines in \texttt{server.log} into \texttt{answer.txt}. \\
\texttt{csv\_sum}         & PC & Add up the \texttt{amount} column of \texttt{sales.csv} and write the total into \texttt{total.txt}. \\
\texttt{find\_errors}     & PC & Copy every line of \texttt{app.log} that contains \texttt{ERROR} into \texttt{errors.log}, in order. \\
\texttt{rename\_photos}   & PC & Rename \texttt{IMG\_1.JPG} \ldots\ \texttt{IMG\_3.JPG} to \texttt{photo\_001.jpg} \ldots\ \texttt{photo\_003.jpg}. \\
\texttt{delete\_tmp}      & PC & Delete every \texttt{.tmp} file in the folder and its subfolders, keeping all other files. \\
\texttt{merge\_parts}     & PC & Join \texttt{part1.txt}, \texttt{part2.txt} and \texttt{part3.txt}, in order, into \texttt{full.txt}. \\
\texttt{scaffold}         & PC & Create a folder \texttt{src} with an empty \texttt{\_\_init\_\_.py}, an empty folder \texttt{tests} and a \texttt{README.md}. \\
\midrule
\multicolumn{3}{l}{\emph{Held-out tasks}} \\
\texttt{reverse\_words}   & coding & Write \texttt{reverse\_words(s)} so the tests pass, without changing the tests. \\
\texttt{average\_empty}   & coding & Fix \texttt{stats.py} so the failing tests pass, without changing the tests. \\
\texttt{failed\_rows}     & PC & Count the rows of \texttt{jobs.csv} whose status is \texttt{failed} and write the number into a file. \\
\texttt{largest\_file}    & PC & Find the largest of four files by size in bytes and write its name into a file. \\
\texttt{us\_spelling}     & PC & Replace every \texttt{colour} with \texttt{color} in all \texttt{.md} files, including subfolders. \\
\bottomrule
\end{tabular}
\caption{The 20 tasks. Each is checked automatically from the final state of its folder.}
\label{tab:tasks}
\end{table}

\section{Results per Task}
\label{app:pertask}

\begin{table}[H]
\centering
\scriptsize
\setlength{\tabcolsep}{3.5pt}
\begin{tabular}{l|ccccccccc|ccc}
\toprule
 & \multicolumn{9}{c|}{qwen 0.8B} & \multicolumn{3}{c}{gemma e2b} \\
Task & C0 & C1 & refiner & gate & memdb & verified & errors & C8 & Reflexion & C0 & C1 & C8 \\
\midrule
\texttt{fix\_fib}         & 1 & 0 & 0 & 0 & 0 & 0 & 0 & 0 & 0 & 3 & 3 & 3 \\
\texttt{palindrome}       & 1 & 1 & 1 & 1 & 0 & 0 & 1 & 0 & 0 & 1 & 2 & 1 \\
\texttt{fix\_import}      & 0 & 0 & 0 & 0 & 0 & 0 & 0 & 0 & 0 & 3 & 2 & 2 \\
\texttt{word\_count}      & 0 & 0 & 1 & 0 & 0 & 0 & 0 & 1 & 0 & 3 & 2 & 3 \\
\texttt{safe\_divide}     & 3 & 0 & 2 & 3 & 1 & 0 & 1 & 2 & 1 & 3 & 3 & 3 \\
\texttt{rename\_function} & 1 & 0 & 0 & 0 & 1 & 1 & 0 & 0 & 0 & 3 & 3 & 3 \\
\texttt{json\_port}       & 3 & 2 & 3 & 2 & 2 & 1 & 1 & 2 & 1 & 3 & 3 & 2 \\
\texttt{organize}         & 0 & 0 & 0 & 0 & 0 & 0 & 0 & 0 & 0 & 3 & 3 & 3 \\
\texttt{count\_lines}     & 2 & 1 & 2 & 0 & 0 & 2 & 0 & 1 & 0 & 3 & 3 & 2 \\
\texttt{csv\_sum}         & 1 & 1 & 1 & 0 & 1 & 0 & 0 & 1 & 0 & 3 & 2 & 2 \\
\texttt{find\_errors}     & 1 & 1 & 1 & 1 & 0 & 0 & 0 & 2 & 1 & 3 & 2 & 3 \\
\texttt{rename\_photos}   & 0 & 0 & 0 & 0 & 0 & 1 & 0 & 1 & 0 & 3 & 2 & 3 \\
\texttt{delete\_tmp}      & 1 & 0 & 0 & 1 & 2 & 1 & 1 & 0 & 0 & 3 & 3 & 3 \\
\texttt{merge\_parts}     & 0 & 1 & 0 & 0 & 1 & 1 & 0 & 1 & 0 & 3 & 3 & 2 \\
\texttt{scaffold}         & 1 & 0 & 0 & 0 & 1 & 1 & 0 & 1 & 0 & 3 & 1 & 2 \\
\midrule
\texttt{reverse\_words}   & 1 & 0 & 1 & 1 & 1 & 0 & 0 & 3 & 0 & 3 & 2 & 2 \\
\texttt{average\_empty}   & 1 & 0 & 0 & 1 & 1 & 0 & 0 & 1 & 0 & 3 & 2 & 3 \\
\texttt{failed\_rows}     & 0 & 1 & 1 & 0 & 0 & 1 & 0 & 2 & 0 & 2 & 1 & 3 \\
\texttt{largest\_file}    & 1 & 1 & 0 & 1 & 1 & 1 & 1 & 1 & 0 & 2 & 2 & 2 \\
\texttt{us\_spelling}     & 1 & 0 & 0 & 1 & 0 & 0 & 0 & 2 & 0 & 2 & 2 & 1 \\
\midrule
Total (of 60)             & 19 & 9 & 13 & 12 & 12 & 10 & 5 & 21 & 3 & 55 & 46 & 48 \\
\bottomrule
\end{tabular}
\caption{Number of seeds (out of 3) in which each task passed at evaluation. The last five rows are
the held-out tasks.}
\label{tab:pertask}
\end{table}

\section{Structured Error Example}
\label{app:errors}

A failed command in C1 returns only the raw output. With structured errors, the same failure
returns a short header first:
\begin{verbatim}
ERROR
type: python_exception
message: ImportError: cannot import name 'add_numbers' from 'mathutils'
hint: The Python code raised an exception. Read the message and the line
      it points to, then fix it.
--- full output ---
exit code: 1
stderr:
Traceback (most recent call last):
  File "main.py", line 1, in <module>
    from mathutils import add_numbers
ImportError: cannot import name 'add_numbers' from 'mathutils'
\end{verbatim}
The 14 error types are:
\begin{center}
\small
\begin{tabular}{lll}
\texttt{outside\_task\_folder} & \texttt{file\_not\_found} & \texttt{blocked\_by\_sandbox} \\
\texttt{timeout} & \texttt{task\_folder\_deleted} & \texttt{bad\_tool\_call} \\
\texttt{unknown\_skill\_or\_subagent} & \texttt{module\_not\_found} & \texttt{command\_not\_found} \\
\texttt{syntax\_error} & \texttt{test\_failed} & \texttt{python\_exception} \\
\texttt{write\_failed} & \texttt{command\_failed} & \\
\end{tabular}
\end{center}

\section{Reflexion Prompt}
\label{app:reflexion}

After each failed train task, the worker receives this prompt, with the task and its own steps
filled in:
\begin{verbatim}
You are an AI agent that does computer tasks (coding, files, shell)
inside a task folder. You just FAILED a task. Below are the task and
your steps.

## Task
{instruction}

## Your steps
{steps}

In a few sentences, diagnose a possible reason for the failure and
write a new, concise, high-level plan that avoids the same mistake.
Write it so it also helps on other, similar tasks.
Reply with the reflection only.
\end{verbatim}

\section{All Settings}
\label{app:settings}

\begin{table}[H]
\centering
\small
\begin{tabular}{lL{9.6cm}}
\toprule
Setting & Value \\
\midrule
Max steps per task & 15 \\
Temperature & 0.7 \\
Context size & 32{,}768 tokens \\
Max output tokens per reply & 4{,}096 \\
Seeds & 3 (0, 1, 2) \\
Sandbox command time limit & 30 seconds \\
Refiner schedule & first after step 25; every 25 steps to step 200; then every 100 \\
Refiner window & last 25 steps (20 for the memory pass) \\
Strategy limit & 4{,}000 characters \\
Memory limit (C1) & 20 notes of 500 characters \\
Skill limit & 10 skills of 3{,}000 characters of code \\
Sub-agent limit & 5 sub-agents, 8 turns each \\
Regression gate & rerun up to 3 recent train tasks with the new harness; keep the edit if it passes at least as many as the old harness did (a tie keeps the edit) \\
Memory database & up to 200 notes; top 5 by cosine similarity (\texttt{nomic-embed-text}) \\
Reflexion & reflect after failed train tasks; keep last 3; 500 characters each \\
Separate refiner & \texttt{gemma4:e4b} \\
\bottomrule
\end{tabular}
\caption{Settings used in all runs.}
\label{tab:settings}
\end{table}

\section{Data and Reproduction}
\label{app:repro}

\paragraph{Code and data.}
The code, tasks, checkers and data are available at
\url{https://github.com/asadsehto/continual-harness-small-models}. The folder \texttt{data/} holds
the outcome of every task run used in this paper (learning, gate reruns and evaluation: 1{,}256
runs), with its pass/fail result, steps, stop reason, tokens and time; every refiner run with the
size of the harness after it and the gate's decision; the final harness size of every seed; skill and
sub-agent calls; and the Ollama version and model file digest of every model. The data contain no
trajectories, so no model output or file content is included.

\paragraph{Rebuilding the tables.}
One command rebuilds every table and chart number in this paper from \texttt{data/}, including the
Holm-adjusted $p$-values, the per-seed results and the determinism check:
\begin{verbatim}
uv run python -m scripts.make_tables
\end{verbatim}

\paragraph{Running the experiments again.}
These commands run the conditions from scratch (they need Ollama and the models):
\begin{verbatim}
uv run python -m scripts.check_tasks        # test every checker
uv run python -m scripts.run_pilot \
    --model <model> --tasks all --seeds 3   # C0
uv run python -m scripts.run_phase1 \
    --model <model> --seeds 3               # C1
uv run python -m scripts.run_phase1 \
    --model <model> --seeds 3 \
    --refiner-model gemma4:e4b --gate --memory-db \
    --verified-skills --structured-errors   # C8
uv run python -m scripts.run_phase1 \
    --model <model> --seeds 3 --reflexion   # Reflexion
uv run python -m scripts.export_data        # runs/ to data/
\end{verbatim}
Each run saves its settings, the server version and model digest, every step, every harness edit
(before and after) and every result. Runs on other hardware or Ollama versions will not repeat our
numbers exactly (Section~\ref{sec:setup}), but should show the same pattern.
